\capitulo{3}{Conceptos teóricos}

\section{Modelo Entidad-Relación}

El modelo Entidad-Relación es un modelo conceptual utilizado para representar la información de un dominio antes de llevarla a una base de datos concreta. Fue propuesto por Chen como una forma de describir la información mediante entidades y relaciones, separando el diseño conceptual de los detalles propios de la implementación final \cite{chen1976}.

Esta separación resulta útil porque permite razonar primero sobre el problema que se quiere representar, sin empezar directamente por tablas, columnas o sentencias SQL. De esta forma, el diseño puede centrarse en identificar qué elementos intervienen, qué información se almacena sobre ellos y cómo se relacionan entre sí.

En el contexto de este trabajo, el modelo Entidad-Relación sirve como base teórica para el editor desarrollado. La aplicación permite representar visualmente diagramas E/R y, a partir de esa representación, aplicar validaciones y transformaciones posteriores hacia el modelo relacional. Por ello, los elementos gráficos del editor no deben entenderse solo como dibujos, sino como partes de un modelo con significado propio.

\section{Elementos principales del modelo}

Un diagrama Entidad-Relación se construye a partir de varios elementos que permiten describir la información de un dominio de forma visual. Aunque existen distintas notaciones y variantes, en este trabajo se consideran especialmente los conceptos que son necesarios para entender el funcionamiento del editor y las transformaciones que se realizan posteriormente.

Los elementos básicos del modelo son los siguientes:

\begin{itemize}
    \item \textbf{Entidades}: representan objetos o conceptos del dominio sobre los que se quiere almacenar información. Por ejemplo, en un sistema académico podrían existir entidades como alumno, asignatura o profesor.

    \item \textbf{Atributos}: describen propiedades de una entidad o de una interrelación. Algunos atributos pueden actuar como identificadores, mientras que otros simplemente aportan información adicional sobre el elemento al que pertenecen.

    \item \textbf{Interrelaciones}: representan asociaciones entre entidades. Estas interrelaciones permiten expresar cómo participan unas entidades con otras dentro del dominio que se está modelando.

    \item \textbf{Cardinalidades}: indican la forma en la que las entidades participan en una interrelación. Son importantes porque condicionan la interpretación del diagrama y también su transformación posterior al modelo relacional.
\end{itemize}

Además de estos elementos básicos, el modelo Entidad-Relación puede incluir otros conceptos más específicos, como entidades débiles, atributos compuestos, atributos multivaluados o interrelaciones de mayor grado. Estos casos permiten representar situaciones más complejas, aunque también hacen necesario mantener una mayor coherencia entre el diagrama, las reglas de validación y la transformación final.

\section{Transformación al modelo relacional}

El modelo Entidad-Relación permite describir un dominio desde un punto de vista conceptual, pero para poder implementarlo en una base de datos relacional es necesario transformarlo a una estructura formada por tablas, columnas y relaciones entre ellas. Esta transformación sirve como paso intermedio entre el diseño conceptual y la generación final de sentencias SQL.

De forma general, las entidades del diagrama suelen convertirse en tablas, sus atributos pasan a ser columnas y los atributos identificadores se utilizan para definir claves primarias. Las interrelaciones entre entidades condicionan la forma en la que se crean claves foráneas o tablas intermedias, dependiendo de la cardinalidad y de la participación de cada entidad en la relación.

Esta fase es importante en el contexto del proyecto porque el editor no se limita a permitir que el usuario dibuje un diagrama. El modelo construido debe poder interpretarse después para obtener una representación relacional coherente. Por ello, la aplicación necesita mantener suficiente información interna sobre entidades, atributos, claves e interrelaciones para poder realizar esta transformación de forma razonable.

\section{Generación de SQL}

SQL es el lenguaje utilizado habitualmente para definir y manipular bases de datos relacionales. Dentro del contexto de este trabajo, interesa especialmente la parte relacionada con la definición de estructuras, como la creación de tablas, claves primarias y claves foráneas.

Una vez que un diagrama Entidad-Relación se ha transformado en un modelo relacional, es posible generar sentencias SQL que representen esa estructura. En este proceso, las tablas resultantes se convierten en sentencias \texttt{CREATE TABLE}, los atributos se traducen a columnas y las relaciones entre tablas se expresan mediante claves foráneas.

En el editor desarrollado, la generación de SQL se entiende como una fase posterior al diseño y a la validación del diagrama. Por tanto, el resultado generado depende de que el modelo interno conserve correctamente la información necesaria sobre entidades, atributos, claves e interrelaciones. De esta forma, el SQL no se genera únicamente a partir de la parte visual del diagrama, sino a partir de la interpretación interna que la aplicación mantiene de él.